
\documentclass{article}
\usepackage{colortbl}
\usepackage{makecell}
\usepackage{multirow}
\usepackage{supertabular}

\begin{document}

\newcounter{utterance}

\centering \large Interaction Transcript for game `hot\_air\_balloon', experiment `air\_balloon\_survival\_en\_reasoning off\_hard', episode 3 with reka{-}flash{-}3.1{-}t0.0.
\vspace{24pt}

{ \footnotesize  \setcounter{utterance}{1}
\setlength{\tabcolsep}{0pt}
\begin{supertabular}{c@{$\;$}|p{.15\linewidth}@{}p{.15\linewidth}p{.15\linewidth}p{.15\linewidth}p{.15\linewidth}p{.15\linewidth}}
    \# & $\;$A & \multicolumn{4}{c}{Game Master} & $\;\:$B\\
    \hline

    \theutterance \stepcounter{utterance}  
    & & \multicolumn{4}{p{0.6\linewidth}}{
        \cellcolor[rgb]{0.9,0.9,0.9}{
            \makecell[{{p{\linewidth}}}]{
                \texttt{\tiny{[P1$\langle$GM]}}
                \texttt{You are participating in a collaborative negotiation game.} \\
\\ 
\texttt{Together with another participant, you must agree on a single set of items that will be kept. Each of you has your own view of how much each item matters to you (importance). You do not know how the other participant values the items. Additionally, you are given the effort each item demands.} \\
\texttt{You may only agree on a set if the total effort of the selected items does not exceed a shared limit:} \\
\\ 
\texttt{LIMIT = 3629} \\
\\ 
\texttt{Here are the individual item effort values:} \\
\\ 
\texttt{Item effort = \{"A21": 5, "B93": 793, "A11": 21, "B33": 842, "C08": 313, "B39": 842, "B96": 631, "C85": 226, "A82": 87, "B31": 761, "B20": 230, "A60": 287, "B09": 698, "C81": 641, "B21": 882\}} \\
\\ 
\texttt{Here is your personal view on the importance of each item:} \\
\\ 
\texttt{Item importance values = \{"A21": 138, "B93": 583, "A11": 868, "B33": 822, "C08": 783, "B39": 65, "B96": 262, "C85": 121, "A82": 508, "B31": 780, "B20": 461, "A60": 484, "B09": 668, "C81": 389, "B21": 808\}} \\
\\ 
\texttt{Goal:} \\
\\ 
\texttt{Your goal is to negotiate a shared set of items that benefits you as much as possible (i.e., maximizes total importance to YOU), while staying within the LIMIT. You are not required to make a PROPOSAL in every message {-} you can simply negotiate as well. All tactics are allowed!} \\
\\ 
\texttt{Interaction Protocol:} \\
\\ 
\texttt{You may only use the following structured formats in a message:} \\
\\ 
\texttt{PROPOSAL: \{'A', 'B', 'C', …\}} \\
\texttt{Propose keeping exactly those items.} \\
\texttt{REFUSE: \{'A', 'B', 'C', …\}} \\
\texttt{Explicitly reject opponent's proposal.} \\
\texttt{ARGUMENT: \{'...'\}} \\
\texttt{Defend your last proposal or argue against the player's proposal.} \\
\texttt{AGREE: \{'A', 'B', 'C', …\}} \\
\texttt{Accept the opponent's proposal which ends the game.} \\
\\ 
\\ 
\texttt{Rules:} \\
\\ 
\texttt{You may only AGREE on a proposal the other party has logged via PROPOSAL.} \\
\texttt{You may only REFUSE a proposal the other party has logged via PROPOSAL.} \\
\texttt{Total effort of any PROPOSAL or AGREE set must be ≤ LIMIT.} \\
\texttt{Do NOT reveal your hidden importance scores.} \\
\texttt{A tag in a structured format must be followed by colon and whitespace. The argument must be a python set containing 0 or more strings.} \\
\texttt{So, it must be of the form TAG: \{...\}} \\
\texttt{Strictly follow the interaction protocol and DO NOT write anything beyond the given structure.} \\
\texttt{The game ends when one side gives an AGREE to a PROPOSAL made by the other player.} \\
\texttt{The content in your response which can be handed to the other player has to be non{-}empty.} \\
\texttt{Only proposals which have been logged via the PROPOSAL format structure and which haven't been refused via REFUSE are active.} \\
\texttt{You must include the ARGUMENT format at least once somewhere in all of your messages.} \\
            }
        }
    }
    & \\ \\

    \theutterance \stepcounter{utterance}  
    & \multicolumn{4}{p{0.6\linewidth}}{
        \cellcolor[rgb]{0.9,0.9,0.9}{
            \makecell[{{p{\linewidth}}}]{
                \texttt{\tiny{[P1$\rangle$GM]}}
                \texttt{<reasoning>} \\
\texttt{Let's tackle this negotiation problem. The goal is to maximize my total importance while keeping the total effort under 3629. First, I need to figure out which items I value the most and which have the least effort. Since the other person's importance isn't known, I have to focus solely on my own importance scores and the effort constraints.} \\
\\ 
\texttt{Let me start by listing all the items with their effort and my importance. Maybe sorting them by importance per unit effort could help prioritize. Wait, but since the effort is a hard limit, maybe I should look for the highest importance items first, but also considering their effort. Alternatively, maybe a knapsack approach where I select items with the highest value (importance) per unit weight (effort) to maximize the total importance without exceeding the limit.} \\
\\ 
\texttt{Let me list out each item with effort and importance:} \\
\\ 
\texttt{Item | Effort | Importance} \\
\texttt{{-}{-}{-}|{-}{-}{-}|{-}{-}{-}} \\
\texttt{A21 | 5 | 138} \\
\texttt{B93 | 793 | 583} \\
\texttt{A11 | 21 | 868} \\
\texttt{B33 | 842 | 822} \\
\texttt{C08 | 313 | 783} \\
\texttt{B39 | 842 | 65} \\
\texttt{B96 | 631 | 262} \\
\texttt{C85 | 226 | 121} \\
\texttt{A82 | 87 | 508} \\
\texttt{B31 | 761 | 780} \\
\texttt{B20 | 230 | 461} \\
\texttt{A60 | 287 | 484} \\
\texttt{B09 | 698 | 668} \\
\texttt{C81 | 641 | 389} \\
\texttt{B21 | 882 | 808} \\
\\ 
\texttt{First, let's compute the importance per effort for each item to see which gives the most bang for the buck. But since the total effort is a limit, maybe a greedy approach by highest importance first, then adjust if needed.} \\
\\ 
\texttt{Alternatively, dynamic programming for the knapsack problem. But since there are 14 items, maybe manageable manually.} \\
\\ 
\texttt{Let me sort the items by importance in descending order:} \\
\\ 
\texttt{1. A11: 868} \\
\texttt{2. B33: 822} \\
\texttt{3. B21: 808} \\
\texttt{4. B09: 668} \\
\texttt{5. A60: 484} \\
\texttt{6. A82: 508 (Wait, A82's importance is 508, which is higher than A60's 484. So maybe I messed up the order. Let me re{-}sort.)} \\
\\ 
\texttt{Wait, let's list all importances:} \\
\\ 
\texttt{A21:138, B} \\
            }
        }
    }
    & & \\ \\

    \theutterance \stepcounter{utterance}  
    & & \multicolumn{4}{p{0.6\linewidth}}{
        \cellcolor[rgb]{0.9,0.9,0.9}{
            \makecell[{{p{\linewidth}}}]{
                \texttt{\tiny{[P1$\langle$GM]}}
                \texttt{Your response did not contain an argument. You must include the structured format ARGUMENT: \{...\} somewhere in your response. Try again} \\
            }
        }
    }
    & \\ \\

    \theutterance \stepcounter{utterance}  
    & \multicolumn{4}{p{0.6\linewidth}}{
        \cellcolor[rgb]{0.9,0.9,0.9}{
            \makecell[{{p{\linewidth}}}]{
                \texttt{\tiny{[P1$\rangle$GM]}}
                \texttt{<reasoning>} \\
\texttt{Let me try to figure out how to approach this. The user mentioned that my previous response didn't include an ARGUMENT. Let me check the rules again. The interaction protocol requires that I must include the ARGUMENT format at least once in my messages. So, in my previous attempt, I probably forgot to add that.} \\
\\ 
\texttt{Alright, let's start over. The goal is to propose a set of items that maximizes my total importance while keeping the total effort under 3629. I need to consider both the effort and my own importance values. The other person's values are unknown, so I can't factor those in. } \\
\\ 
\texttt{First, I should list all items with their effort and my importance. Let me do that:} \\
\\ 
\texttt{Item | Effort | Importance} \\
\texttt{{-}{-}{-}|{-}{-}{-}|{-}{-}{-}} \\
\texttt{A21 | 5 | 138} \\
\texttt{B93 | 793 | 583} \\
\texttt{A11 | 21 | 868} \\
\texttt{B33 | 842 | 822} \\
\texttt{C08 | 313 | 783} \\
\texttt{B39 | 842 | 65} \\
\texttt{B96 | 631 | 262} \\
\texttt{C85 | 226 | 121} \\
\texttt{A82 | 87 | 508} \\
\texttt{B31 | 761 | 780} \\
\texttt{B20 | 230 | 461} \\
\texttt{A60 | 287 | 484} \\
\texttt{B09 | 698 | 668} \\
\texttt{C81 | 641 | 389} \\
\texttt{B21 | 882 | 808} \\
\\ 
\texttt{Now, I need to select a subset of these items such that the sum of their efforts is ≤ 3629, and the sum of my importances is maximized. This is a classic 0/1 knapsack problem. Since the limit is 3629, which is a pretty high number, but there are items with high effort, I need to balance high importance items with their effort.} \\
\\ 
\texttt{Let me sort the items by importance in descending order to see which ones are most valuable to me:} \\
\\ 
\texttt{1. A11 (868)} \\
\texttt{2. B33 (822)} \\
\texttt{3. B21 (808)} \\
\texttt{4. B09 (668)} \\
\texttt{5. A82 (508)} \\
\texttt{6. A60 (484)} \\
\texttt{7. C08 (783) – Wait, C08's importance is 783, which is higher than A82's 508. So maybe I need to re{-}sort properly.} \\
\\ 
\texttt{Wait, let's list all importances:} \\
\\ 
\texttt{A11: 868} \\
\texttt{B33: 822} \\
\texttt{B} \\
            }
        }
    }
    & & \\ \\

    \theutterance \stepcounter{utterance}  
    & & \multicolumn{4}{p{0.6\linewidth}}{
        \cellcolor[rgb]{0.9,0.9,0.9}{
            \makecell[{{p{\linewidth}}}]{
                \texttt{\tiny{[P1$\langle$GM]}}
                \texttt{Your response did not contain an argument. You must include the structured format ARGUMENT: \{...\} somewhere in your response. Try again} \\
            }
        }
    }
    & \\ \\

    \theutterance \stepcounter{utterance}  
    & \multicolumn{4}{p{0.6\linewidth}}{
        \cellcolor[rgb]{0.9,0.9,0.9}{
            \makecell[{{p{\linewidth}}}]{
                \texttt{\tiny{[P1$\rangle$GM]}}
                \texttt{<reasoning>} \\
\texttt{Let me try to figure out how to approach this. The user mentioned that my previous response didn't include an ARGUMENT. Let me check the rules again. The interaction protocol requires that I must include the ARGUMENT format at least once in my messages. So, in my previous attempt, I probably forgot to add that.} \\
\\ 
\texttt{Alright, let's start over. The goal is to propose a set of items that maximizes my total importance while keeping the total effort under 3629. I need to consider both the effort and my own importance values. The other person's values are unknown, so I can't factor those in. } \\
\\ 
\texttt{First, I should list all items with their effort and my importance. Let me do that:} \\
\\ 
\texttt{Item | Effort | Importance} \\
\texttt{{-}{-}{-}|{-}{-}{-}|{-}{-}{-}} \\
\texttt{A21 | 5 | 138} \\
\texttt{B93 | 793 | 583} \\
\texttt{A11 | 21 | 868} \\
\texttt{B33 | 842 | 822} \\
\texttt{C08 | 313 | 783} \\
\texttt{B39 | 842 | 65} \\
\texttt{B96 | 631 | 262} \\
\texttt{C85 | 226 | 121} \\
\texttt{A82 | 87 | 508} \\
\texttt{B31 | 761 | 780} \\
\texttt{B20 | 230 | 461} \\
\texttt{A60 | 287 | 484} \\
\texttt{B09 | 698 | 668} \\
\texttt{C81 | 641 | 389} \\
\texttt{B21 | 882 | 808} \\
\\ 
\texttt{Now, I need to select a subset of these items such that the sum of their efforts is ≤ 3629, and the sum of my importances is maximized. This is a classic 0/1 knapsack problem. Since the limit is 3629, which is a pretty high number, but there are items with high effort, I need to balance high importance items with their effort.} \\
\\ 
\texttt{Let me sort the items by importance in descending order to see which ones are most valuable to me:} \\
\\ 
\texttt{1. A11 (868)} \\
\texttt{2. B33 (822)} \\
\texttt{3. B21 (808)} \\
\texttt{4. B09 (668)} \\
\texttt{5. C08 (783) – Wait, C08's importance is 783, which is higher than A82's 508. So maybe I need to re{-}sort properly.} \\
\\ 
\texttt{Wait, let's list all importances:} \\
\\ 
\texttt{A11: 868} \\
\texttt{B33: 822} \\
\texttt{B21: 808} \\
\texttt{B09: 668} \\
\texttt{C08:} \\
            }
        }
    }
    & & \\ \\

    \theutterance \stepcounter{utterance}  
    & & & \multicolumn{2}{p{0.3\linewidth}}{
        \cellcolor[rgb]{0.9,0.9,0.9}{
            \makecell[{{p{\linewidth}}}]{
                \texttt{\tiny{[GM$|$GM]}}
                \texttt{abort} \\
            }
        }
    }
    & & \\ \\

    \theutterance \stepcounter{utterance}  
    & & & \multicolumn{2}{p{0.3\linewidth}}{
        \cellcolor[rgb]{0.9,0.9,0.9}{
            \makecell[{{p{\linewidth}}}]{
                \texttt{\tiny{[GM$|$GM]}}
                \texttt{end game} \\
            }
        }
    }
    & & \\ \\

\end{supertabular}
}

\end{document}
